\section{Minimal-circuit implementation of the Legendre character}
\label{sec:legendre-channel}

The modular-multiplication algorithm of \cref{sec:mod-mul-algorithm}
implements each key frequency $k \in \mathcal{K}$ via a cluster of MLP
neurons specialized to that frequency.  At the baseline width
$d_{\text{model}} = 128$, those clusters all live on \emph{primitive}
characters (order $p-1 = 112$).  In this note we report a qualitatively
different feature that appears when the model is forced to use less
capacity: at $d_{\text{model}} = 24$, the network groks by
incorporating the order-$2$ \emph{Legendre} character $\chi_{56}$ into
$\mathcal{K}$, and implements that channel with only $9$ MLP neurons --
roughly $10$--$20\times$ fewer than each primitive channel in the same
model.  The implementation is the minimal possible ReLU circuit that
computes $\chi_{56}(ab)$: two thresholded half-planes detecting the two
same-coset blocks of $\Zpstar{p} \times \Zpstar{p}$, summed into a
single residual-stream direction that the unembed reads as the $k=56$
component.  We treat this as evidence that under capacity pressure the
model preserves the algebraic structure of the algorithm of
\cref{subsec:mul-algorithm} but selects the cheapest available
implementation of each individual channel.

\begin{figure}[ht]
  \centering
  \includegraphics[width=0.96\linewidth]{legendre_reconstructed_signal}
  \caption{The punchline.  Top row: the residual-stream signal that the
    $d_{\text{model}} = 24$ model's $9$-neuron Legendre cluster writes
    along its dominant $\Wout$ direction, plotted as a function of the
    input pair $(a, b) \in \Zpstar{p} \times \Zpstar{p}$.  Bottom row: the
    mathematical reference $-\chi_{56}(a)\,\chi_{56}(b) = -\chi_{56}(ab)$.
    Left column: natural integer indexing of $a$ and $b$.  Right
    column: rows and columns reordered by coset (QRs first, then
    non-QRs).  The cluster output reproduces $-\chi_{56}(ab)$
    pointwise; the rank-$1$ residual direction is $0.998$-aligned with
    the unembed's $k = 56$ read direction, so the cluster's contribution
    to the logits is exactly the $k = 56$ term of
    \eqref{eq:logit-formula}.  The rest of this note unpacks how the
    network achieves this with $9$ ReLU neurons rather than the $\sim
    100$ used by each primitive channel.}
  \label{fig:legendre-reconstructed}
\end{figure}

\subsection{Setup}
\label{subsec:legendre-setup}

\Cref{sec:mod-mul-algorithm} pins the baseline model at
$d_{\text{model}} = 128$.  To map the minimum width at which grokking
still occurs we trained a sweep at $d_{\text{model}} \in \{12, 16, 24,
32, 64, 128\}$ holding all other hyperparameters fixed.  Models at
$d_{\text{model}} \geq 32$ grok in the usual $\sim 5\text{--}10\mathrm{k}$
epochs.  The $d_{\text{model}} = 24$ run was carried to $60\mathrm{k}$
epochs; it eventually groks at epoch $\approx 20{,}800$.  The runs at
$d_{\text{model}} \in \{12, 16\}$ do not grok within the budgeted
$50\mathrm{k}$ epochs.

Recovered key frequencies from the embedding spectrum on each groked run:
\begin{center}
\begin{tabular}{r|l|l|c}
  $d_{\text{model}}$ & $\mathcal{K}$ &
    \multicolumn{1}{c|}{order multiset} & final test loss \\
  \hline
  24  & $\{1,\,9,\,23,\,30,\,56\}$ & $\{112,112,112,56,\mathbf{2}\}$
      & $4.30 \times 10^{-7}$ \\
  32  & $\{3,\,4,\,35,\,47\}$      & $\{112,28,16,112\}$
      & $2.17 \times 10^{-7}$ \\
  64  & $\{3,\,9,\,11,\,47\}$      & $\{112,112,112,112\}$
      & $1.94 \times 10^{-7}$ \\
  128 & $\{17,\,18,\,22,\,41\}$    & $\{112,56,56,112\}$
      & $5.15 \times 10^{-7}$ \\
\end{tabular}
\end{center}

The $d_{\text{model}} = 24$ model is the only run in the sweep whose
$\mathcal{K}$ contains a character of order $2$: that is the Legendre
symbol $\chi_{56}(a) = \left(\frac{a}{p}\right)$.  All other groked
runs use four characters of order $\geq 28$.

\subsection{The 9-neuron Legendre cluster}
\label{subsec:legendre-cluster}

Following the per-neuron analysis of \cref{subsec:implementation}, we
assign each MLP neuron a dominant frequency by projecting its
post-activation tensor $a_n \in \R^{p \times p}$ into the
two-dimensional multiplicative-Fourier basis and summing energy over
the four matched bigrams of each frequency $k$.  At
$d_{\text{model}} = 24$ every neuron clusters cleanly onto $\mathcal{K}$
(no off-$\mathcal{K}$ tail):

\begin{center}
\begin{tabular}{r|r|r}
  cluster $k$ & order & \# neurons \\ \hline
   1 & 112 &  89 \\
   9 & 112 & 170 \\
  23 & 112 &  95 \\
  30 & 112 & 149 \\
  \textbf{56} & \textbf{2} & \textbf{9} \\ \hline
  total &  & 512 \\
\end{tabular}
\end{center}

\begin{figure}[ht]
  \centering
  \includegraphics[width=0.95\linewidth]{neuron_freq_distribution}
  \caption{Distribution of MLP neurons by dominant multiplicative-Fourier
    frequency in the $d_{\text{model}} = 24$ model.  All $512$ neurons
    cluster on $\mathcal{K} = \{1, 9, 23, 30, 56\}$; the order-$2$
    Legendre cluster $k=56$ holds $9$ neurons compared to $89$--$170$
    for each primitive cluster.}
  \label{fig:legendre-distribution}
\end{figure}

The Legendre cluster is $\sim 10$--$20\times$ smaller than each of the
four primitive clusters.  Two questions follow.  First, are those $9$
neurons actually load-bearing or are they a vestigial detail?  Second,
do they implement the same Fourier-multiplication channel as the
primitive clusters, or is the model running a structurally different
algorithm at $k = 56$?

\subsection{Ablation: the Legendre cluster is load-bearing}
\label{subsec:legendre-ablation}

We zero each cluster's post-ReLU output via a forward hook on
\verb|blocks.0.mlp.hook_post|, leaving all other weights untouched, and
report train/test cross-entropy under the original $30\%$ train split.
Baseline (no ablation): train $1.10 \times 10^{-7}$, test $4.30
\times 10^{-7}$, accuracy $1.0000$.

\begin{center}
\begin{tabular}{r|r|r|r|r|r}
  cluster $k$ & cluster size & train loss & test loss & accuracy &
    test ratio vs.\ baseline \\ \hline
  1  &  89 & $3.17 \times 10^{-4}$ & $8.14 \times 10^{-4}$ & $0.9999$ & $1.9 \times 10^{3}$ \\
  9  & 170 & $1.32$ & $1.44$ & $0.6094$ & $3.3 \times 10^{6}$ \\
  23 &  95 & $0.80$ & $0.87$ & $0.7129$ & $2.0 \times 10^{6}$ \\
  30 & 149 & $1.52$ & $1.59$ & $0.6117$ & $3.7 \times 10^{6}$ \\
  \textbf{56} &  \textbf{9} & $\bm{0.55}$ & $\bm{0.64}$ & $\bm{0.7440}$ &
    $\bm{1.5 \times 10^{6}}$ \\
\end{tabular}
\end{center}

Two observations.

\emph{The Legendre cluster does as much algorithmic work as a primitive
cluster, despite holding $\sim 1.8\%$ of $d_{\mathrm{mlp}}$.}  Removing
its $9$ neurons crashes test accuracy from $100\%$ to $74.4\%$ -- the
same magnitude of damage as removing $95$--$170$ neurons from a
primitive cluster.  A size-matched control (ablate $9$ random neurons
drawn from any primitive cluster, $20$ trials) leaves test loss within
$1.4\times$--$3.4\times$ of baseline and accuracy at $1.0$; the Legendre
$9$ alone produce a $\sim 10^{6}\times$ test-loss increase.  This is
\emph{not} a generic ``$9$ neurons matter'' effect.

\emph{Not every primitive cluster is load-bearing.}  Removing all $89$
neurons of the $k = 1$ cluster increases test loss to $8.14 \times
10^{-4}$ but leaves accuracy at $0.9999$.  Despite carrying substantial
energy in the embedding spectrum, $k=1$ is not actually needed at
inference -- whereas $k = 56$, with one-tenth the neuron budget, is.
The neuron-side $\mathcal{K}$ is not identical to the embedding-side
$\mathcal{K}$.

The ablation alone is consistent with two distinct hypotheses: the
Legendre cluster could be a standard additive channel in the
Fourier-multiplication algorithm (removing one term from
\cref{eq:logit-formula} introduces $\pm 1$ errors throughout the logit
vector and would crash any cluster), or it could implement a
qualitatively different mechanism such as a parity-gate.  We
distinguish these in the next subsection.

\subsection{Structure of the cluster: two tiles, four-and-five replicas}
\label{subsec:legendre-structure}

\begin{figure}[ht]
  \centering
  \includegraphics[width=0.96\linewidth]{legendre_grid_v2}
  \caption{Post-ReLU activation $a_n(a, b)$ for each of the $9$ Legendre
    neurons.  Top row: the $5$ neurons of Group A (firing pattern
    identical across the row).  Bottom row: the $4$ neurons of Group B
    (also identical across the row).  Each group is a single
    function replicated; the two groups partition $\Zpstar{p} \times \Zpstar{p}$ into
    disjoint firing supports.}
  \label{fig:legendre-grid}
\end{figure}

Pairwise binarized intersection-over-union of fired masks across the
$9$ Legendre neurons gives a strict $\{1.00, 0.02\}$ structure: $5$
neurons (Group A: $\{207, 287, 336, 418, 481\}$) share a single fired
mask, $4$ neurons (Group B: $\{252, 291, 297, 347\}$) share a different
fired mask, and the across-group IoU is essentially zero.  Block-wise
means in the coset partition of $\Zpstar{p} \times \Zpstar{p}$ (rows
sorted by parity of $\dlog_g a$, columns by parity of $\dlog_g b$):

\begin{center}
\begin{tabular}{r|c|c}
  block & Group A neuron $\#207$ & Group B neuron $\#252$ \\ \hline
  QR $\times$ QR              & $0.000$ & $5.301$ \\
  QR $\times$ non-QR          & $0.064$ & $0.074$ \\
  non-QR $\times$ QR          & $0.060$ & $0.068$ \\
  non-QR $\times$ non-QR      & $4.909$ & $0.000$ \\
\end{tabular}
\end{center}

Each group fires only on a \emph{single} same-coset block: Group A
detects $(\chi_{56}(a), \chi_{56}(b)) = (-1, -1)$, Group B detects
$(+1, +1)$.  The two off-diagonal blocks -- where $ab$ is itself a
non-QR -- are at zero (modulo tiny boundary leakage on the rows/columns
where $a = 0$ or $b = 0$, which lie outside $\Zpstar{p}$); the
coset-sorted view in \cref{fig:legendre-reconstructed} (right column)
displays the same $2 \times 2$ block structure once both groups are
summed.

Each of the two groups is one ReLU tile (with $4$ or $5$ exact
replicas).  Within each group all replicas have identical post-ReLU
output up to scalar magnitude.  Two tiles is exactly the number of
threshold-step functions needed to compute a binary-valued character:
one indicator per value in the image.

\subsection{The residual-stream signal: rank-$1$ with the $k = 56$ unembed direction}
\label{subsec:legendre-Wout}

The $9$-by-$d_{\text{model}}$ submatrix $\Wout[:, \text{Legendre}]$
factorizes cleanly.  Its singular values are
\[
  S \;=\; \bigl[\,3.637,\;\, 0.132,\;\, 4.6 \times 10^{-6},\;\, \ldots \bigr]
\]
so the matrix is effectively \emph{rank-$1$}: $99.87\%$ of the Frobenius
energy lies in the top component.  All $9$ neurons project onto the
same residual-stream direction $v \in \R^{24}$, and -- critically --
they all do so with the \emph{same sign} (per-neuron projections all
$\approx -1.17$ for Group A or $\approx -1.27$ for Group B).  In other
words: both groups contribute additively to a single residual-stream
scalar.

Let $\hat u_{56} := \WU \cos(\omega_{56} \dlog_g \cdot)/\|\cdot\|$
denote the unit-norm $k = 56$ readout direction in the unembed.  We
find
\begin{equation}
  \bigl|\langle v, \hat u_{56}\rangle\bigr| \;=\; 0.998,
  \label{eq:legendre-Wout-alignment}
\end{equation}
i.e.\ the cluster writes essentially exactly into the direction that
the unembed projects out as the $k = 56$ component.  Aggregating across
all $9$ neurons, the scalar signal $s(a, b) := \sum_n a_n(a, b)
\,\langle \Wout[:, n], v\rangle$ has block-wise means $\{-26.83,\,
-0.75,\, -0.69,\, -28.68\}$ in the coset partition above.  Its
correlation with the reference patterns is
\begin{equation}
  \mathrm{corr}\bigl(s(a, b),\; \chi(a)\chi(b)\bigr) = -0.9965,
  \qquad
  \mathrm{corr}\bigl(s(a, b),\; \chi(a)+\chi(b)\bigr) = +0.0485.
  \label{eq:legendre-correlation}
\end{equation}
The cluster output is, to within $0.4\%$, an affine function of
$\chi_{56}(ab)$ -- the DC offset is absorbed by bias terms in the
unembed.

(\Cref{fig:legendre-reconstructed} shows $s(a, b)$ alongside the
mathematical reference $-\chi(ab)$.)

\subsection{Why two tiles suffice}
\label{subsec:legendre-algebra}

The mod-$\!2$ structure of $\Zpstar{p}$ supplies an algebraic identity
that turns $\chi_{56}(ab)$ into a sum of two indicators:
\begin{equation}
  \chi_{56}(ab)
  \;=\; -1
   \;+\; 2\,\mathbf{1}\bigl[\chi(a) = \chi(b) = +1\bigr]
   \;+\; 2\,\mathbf{1}\bigl[\chi(a) = \chi(b) = -1\bigr].
  \label{eq:legendre-identity}
\end{equation}
Both indicators on the right-hand side are linear-threshold functions
of $(\chi(a), \chi(b)) \in \{-1, 0, +1\}^2$ and so can each be realized
by a \emph{single} ReLU tile.  The DC offset $-1$ is absorbed into the
unembed bias.  Concretely, the model's circuit is:
\begin{enumerate}[leftmargin=2em]
  \item The embedding $\WE$ places the $k = 56$ component of token $a$
    along a direction in $\R^{d_{\text{model}}}$ with magnitude
    $\propto \chi_{56}(a)$, and likewise for $b$.
  \item Attention averages the $a$- and $b$-token residuals at the
    ``$=$'' position, so the MLP receives an input whose $k=56$
    component is $\propto \chi_{56}(a) + \chi_{56}(b) \in \{-2, 0, +2\}$.
  \item Group~A neurons fire as $\mathrm{ReLU}(-\chi(a) - \chi(b) -
    \theta)$ with $\theta \in (1, 2)$; they activate uniquely on
    $\chi(a) = \chi(b) = -1$.  Group~B neurons fire as
    $\mathrm{ReLU}(\chi(a) + \chi(b) - \theta)$; they activate
    uniquely on $\chi(a) = \chi(b) = +1$.
  \item All $9$ neurons project their post-ReLU output onto a common
    residual-stream direction $v$ with the same sign of $\Wout$
    coefficient.  By
    \cref{eq:legendre-identity},
    the sum $\sum_n \mathrm{post}_n \cdot \langle \Wout[:, n], v\rangle$
    is proportional to $-\chi_{56}(ab)$ plus a constant.
  \item The unembed reads $v$ via
    $\cref{eq:legendre-Wout-alignment}$, multiplies by $\chi_{56}(c)$
    for each candidate $c$, and contributes $\chi_{56}(c)\,\chi_{56}(ab)
    = \cos(\omega_{56}(\dlog_g ab - \dlog_g c))$ to the logit.
\end{enumerate}

This is the same Fourier-multiplication algorithm of
\cref{subsec:mul-algorithm}.  What is new is the
\emph{circuit instantiation} of one channel.  A primitive character
$\chi_k$ takes $112$ distinct values, and the ReLU MLP approximates
$\cos(\omega_k(\dlog a + \dlog b))$ by tiling those phases -- empirically
$\sim 100$ neurons per primitive channel at $d_{\text{mlp}} = 512$.  An
order-$2$ character takes only two distinct values, so the cosine
\emph{is} a step function and the entire channel reduces to a pair of
linear-threshold detectors. The number of replicas of each tile
($4$ and $5$) reflects gradient-descent's redundancy budget rather
than additional computational structure.

\subsection{Discussion}
\label{subsec:legendre-discussion}

Two themes connect to the rest of the project.

\emph{Capacity pressure selects character orders.}  The cheapest
character to implement, in ReLU MLP units, is the unique order-$2$
character (one tile per value, $|\text{image}| = 2$).  The next
cheapest is an order-$4$ character ($|\text{image}| = 4$, requiring
$\sim 4$ tiles), and so on; primitive characters at order $p - 1$ are
the most expensive.  A model with $d_{\text{model}} = 24$ has roughly
$1/5$ of the baseline's $d_{\text{model}}$, and the embedding's $k =
56$ direction costs half the parameters of a primitive direction
(only a $\cos$ row, since $\sin(\pi \dlog_g a) = 0$ identically).  The
model's spontaneous incorporation of $\chi_{56}$ into $\mathcal{K}$ is
consistent with a cost-driven preference for cheap characters under
tight capacity.

\emph{The model reaches for a cryptographer's primitive.}  The
Legendre symbol $\left(\frac{a}{p}\right)$ is one of the most-studied
objects in computational number theory; constant-time evaluation is the
subject of, e.g., the divsteps-based Kronecker-symbol algorithm of
Aranha et al.\ \cite{aranha2024kronecker}.  That a capacity-limited
mech-interp model spontaneously builds a clean Legendre-channel circuit
-- and, indeed, builds the structurally minimal such circuit -- is a
small but concrete instance of the program of comparing ML-learned
algorithms against cryptographer-designed primitives.  The
$d_{\text{model}} = 24$ network is a tiny existence proof that the
``ML reaches for the Legendre symbol'' phenomenon can be made
mechanistically explicit at a level of detail comparable to the
Fourier-multiplication algorithm itself.
